%======================================================%
%   Beamer Presentation
%   LaTeX Template
%   compile using PDFTeXify or PDFLaTeX
%======================================================%

%--------------------------------------------------------------------------------
%   PACKAGES AND THEMES
%--------------------------------------------------------------------------------

%\documentclass[notheorems,11pt,compress]{beamer}
\documentclass[notheorems,11pt]{beamer}

\usetheme{moloch}
\usepackage{appendixnumberbeamer}

\usepackage{booktabs}
\usepackage[scale=2]{ccicons}

\usepackage{pgfplots}
\usepgfplotslibrary{dateplot}

\usepackage[english]{babel}  % Show english numerate
\molochset{block=fill}
\usepackage{xspace}
\newcommand{\themename}{\textbf{\textsc{metropolis}}\xspace}

%\usefonttheme[onlymath]{serif}
\usefonttheme{serif}
%\setbeamercovered{transparent}
%\usetheme{boxes}

%\setbeamertemplate{headline}{}
%\setbeamertemplate{footline}{}
%\setbeamertemplate{blocks}[rounded][shadow=true]
%\setbeamertemplate{blocks}[default]
\setbeamertemplate{blocks}[shadow=true]
%\setbeamercolor{block body}{bg=black!2}
\setbeamertemplate{navigation symbols}{}
%\setbeamertemplate{itemize items}[square]  % ball, circle
%\setbeamertemplate{enumerate items}[square]
%\setbeamertemplate{section in toc}[square]

\definecolor{iron}{RGB}{0,82,67}
\definecolor{footed}{rgb}{0.5,0.2,0.5}
\setbeamercolor{frametitle}{bg=iron}
%\setbeamercolor{progress bar}{fg=iron,bg=iron}

%分式使用行间版
\everymath{\displaystyle}

\setbeamertemplate{footline}
{%
\leavevmode%
%\hskip1cm \fangsong \insertshortauthor~(\insertshortinstitute) ~~\textbar~~ \insertshorttitle
\hfill\normalfont\color{footed}\insertframenumber\,/\,\inserttotalframenumber \kern1em \hfill
\vskip4pt%
}


%----------- 宏包 -------------
%\usepackage[latin1]{inputenc}
%\usepackage{times}
%\usepackage[T1]{fontenc}

\usepackage[UTF8,noindent]{ctex}
%\usepackage[english]{babel}
\usepackage{amsmath,amssymb,version}
\usepackage{graphicx,fancybox,mathrsfs,multirow}
\usepackage{booktabs}
\usepackage{epsfig,epstopdf}
\usepackage{url,hyperref}
\usepackage{tabularx,array,makecell}
\usepackage{color,xcolor}
\usepackage{cases}
\usepackage{mathtools}
\usepackage{tikz}

\usepackage{scrextend}

%---------- 定义表格命令 ----------
\newcolumntype{P}[1]{>{\centering \arraybackslash}p{#1}}
\newcolumntype{L}{X}
\newcolumntype{C}{>{\centering \arraybackslash}X}
\newcolumntype{R}{>{\raggedleft \arraybackslash}X}



%---------- 设置 定理环境  --------------
\theoremstyle{plain}
\setbeamertemplate{theorems}[numbered]
%\theoremstyle{definition}
%\newtheorem{theorem}{Theorem}
\newtheorem{theorem}{定理}
\numberwithin{theorem}{section}
%\newtheorem{definition}{Definition}
\newtheorem{definition}{定义}
\numberwithin{definition}{section}
%\newtheorem{lemma}{Lemma}
\newtheorem{lemma}{引理}
\numberwithin{lemma}{section}
%\newtheorem{proposition}{Proposition}
\newtheorem{proposition}{命题}
\numberwithin{proposition}{section}
%\newtheorem{corollary}{Corollary}
\newtheorem{corollary}{{推论}}
\numberwithin{corollary}{section}
\theoremstyle{example}
%\newtheorem{example}{Example}
\newtheorem{example}{例}
%\numberwithin{example}{section}
\renewenvironment{proof}[1][证明]{\textbf{#1}:~}{\qed\par}
%\renewenvironment{proof}[1][证明]{{\heiti #1}:~}{\qed\par}


\setbeamerfont{caption name}{family=\rmfamily,series=\normalfont}
\setbeamertemplate{caption}[numbered]
\numberwithin{figure}{section}
\numberwithin{table}{section}
%\numberwithin{equation}{section}

%\addtokomafont{labelinglabel}{\normalfont}%\bfseries
%\usepackage{xunicode-addon}
%\renewfontfamily\useTIPAfont{Segoe UI}

%---------- 设置字体  ---------------
%\setbeamerfont{normal text}{family=\rmfamily}
%\setbeamerfont{frametitle}{family=\sffamily}
%\setbeamerfont{title}{family=\sffamily}
%\setbeamerfont{subtitle}{family=\sffamily}
%\setbeamerfont{institute}{family=\rmfamily}
%\setbeamerfont{author}{family=\rmfamily}
%\setbeamerfont{date}{family=\rmfamily}
%\setbeamerfont{headline}{family=\sffamily}
%\setbeamerfont{footline}{family=\rmfamily}
%\setbeamerfont{section in toc}{family=\rmfamily}
%\setbeamerfont{subsection in toc}{family=\rmfamily}
%\AtBeginDocument{\usebeamerfont{normal text}}


%---------- 定义行距 ----------
\renewcommand{\baselinestretch}{1.15}

%---------- 调节公式的间距 ----------
%\AtBeginDocument{
%	\setlength{\abovedisplayskip}{4pt plus 1pt minus 1pt}
%	\setlength{\belowdisplayskip}{4pt plus 1pt minus 1pt}
%	\setlength{\abovedisplayshortskip}{2pt}
%	\setlength{\belowdisplayshortskip}{2pt}
%	\setlength{\arraycolsep}{2pt}
%}

\makeatletter
\newcommand\HUGE{\@setfontsize\Huge{36}{42}}
\makeatother

%---------- 自定义命令 ----------
\newcommand{\red}[1]{\textcolor{red}{#1}}
\newcommand{\blue}[1]{\textcolor{blue}{#1}}


%\AtBeginSection[]{
%\begin{frame}
%  \frametitle{目录}
%  \tableofcontents[currentsection,currentsubsection]
%  %\tableofcontents[currentsection,currentsubsection,subsectionstyle=show/show/shaded]
%  \addtocounter{framenumber}{-1}
%\end{frame}
%}


%\setmonofont{Consolas}
%\setmonofont{Latin Modern Mono}
%\title{Metropolis}
%\subtitle{A modern beamer theme}
%\date{\today}
%\author{Matthias Vogelgesang}
%\institute{Center for modern beamer themes}

%----------------------------------------------------------------------------------------
%	TITLE PAGE
%----------------------------------------------------------------------------------------

\title[Short 题目]{运动规划中的混合整数规划} % The short title appears at the bottom of every slide, the full title is only on the title page

\author{姓名}  % 名字
\institute[NU] % 机构缩写
{
Name of University \\ % 机构全称
\medskip
\textit{name@email.com} % 邮件
}
\date[2020.6.23]{2020~年~6~月~23~日} % 日期


\graphicspath{{./figures/}}

\begin{document}

% only change the spacing of body text
%\setlength{\baselineskip}{15pt}

\begin{frame}
\titlepage
\end{frame}

\begin{frame}{目录}
  \setbeamertemplate{section in toc}[sections numbered]
\tableofcontents
\end{frame}

%----------------------------------------------------------------------------------------
%	PRESENTATION SLIDES
%----------------------------------------------------------------------------------------

%------------------------------------------------
\section{问题背景}
%------------------------------------------------

\begin{frame}{混合整数规划}
\begin{block}{混合整数规划}
混合整数规划(Mixed Integer Programming, MIP)是一类数学规划问题，其目标是优化一个线性函数，其中部分决策变量为整数.
\end{block}
\begin{enumerate}[<+-| alert@+>]
  \item 涉及整数的离散问题 \par
  \begin{itemize}[<+-| alert@+>]
    \item 背包问题
  \end{itemize}
  \item 涉及逻辑条件的问题
  \begin{itemize}[<+-| alert@+>]
    \item 复杂运输过程中的碰撞避免
  \end{itemize}
  \item 排序和分配问题
  \begin{itemize}[<+-| alert@+>]
  \item 任务分配
  \item 旅行商问题
\end{itemize}
\end{enumerate}

\end{frame}

\begin{frame}{MIP的特点与优势}
\begin{itemize}[<+-| alert@+>]
  \item 可以非线性建模
  \item 允许非凸约束
  \item 可以包含逻辑决策
  \item 已广泛应用于控制工程问题
\end{itemize}
\end{frame}

\begin{frame}{符号与记号}
  \begin{enumerate}
    \item 逻辑运算符：$\wedge$(与), $\vee$(或), $\neg$(非)
    \item 集合的Minkowski和：$A\oplus B=\{x:x=a+b,a\in A,b\in B\}$
    \item $\|x\|^2_Q=x^TQx$
    \item 补集$C_X(S)$,闭包$\operatorname{cl}(S)$
    \item $\mathcal{V}(P)$ 表示多面体$P\subset \mathbb{R}^n$ 的(有限)顶点集
    \item $\#\mathcal{I}$表示离散集$\mathcal{I}$的基数
  \end{enumerate}
  \vspace{1 em}
  
\end{frame}

%-----------------------------------------------
\section{MIP公式化}

\subsection{广义析取规划}
%-----------------------------------------------
\begin{frame}{广义析取规划(GDP)}
  \begin{exampleblock}{GDP的例子(1)}
    \begin{subequations}
    \begin{align}
      &\min_x{f(x)+\sum_k}{c_k} \\
      \text{s.t. } &r(x)\leqslant 0,x\in\mathbb{R}^n,c_k\in\mathbb{R} \\
      &\bigvee_{j\in J_k}\left[\begin{gathered}
        Y_{jk} \\
        g_{jk}(x)\leqslant 0 \\
        c_k=\gamma_{jk}
      \end{gathered}\right],k\in K \label{1c}\\
      &\Omega(Y)=\text{true},Y_{jk}\in\{\text{true,false}\} \label{1d}
    \end{align}
    \end{subequations}
  \end{exampleblock}
\end{frame}

\begin{frame}{问题1的MIP改写}
使用二元变量$y_{jk}=\{0,1\}$代替布尔变量$Y_{jk}$，并将约束(\ref{1c})替换为：
\begin{subequations}
\begin{align}
  &g_{jk}(x)\leqslant M_{jk}(1-y_{jk}) \\
  &\sum_{j\in J_k}{y_{jk}}=1,k\in K
\end{align}
\end{subequations}
其中$M_{jk}$ 是“大M”参数，为足够大的常数.\par
成本变量$c_k$看可以被重写为乘积$\gamma_{jk}y_{jk}$ ，同样地，(\ref{1d})可以重写为线性约束
$$Ay\leqslant a$$
\end{frame}

\begin{frame}{“大M”参数的选取}
  \begin{example}
    原逻辑命题为$x>0\Longrightarrow \alpha=1$，其中$x\in\mathbb{X}=\{x:|x|\leqslant \overline{x}\}$.\par
    引入大M参数后，可将命题转化为不等式$x-M\alpha\leqslant 0$.\par
    为了使不等式满足命题的逻辑含义，需选取$M\geqslant \overline{x}$.\par
    因此最后取$M=\max_{x\in\mathbb{X}}{x}$，即$M=\overline{x}$.
  \end{example}

\end{frame}


\section{几何视角下的MIP}

\begin{frame}{多面体的半空间表示}
  \begin{proposition}{}
    有界多面体$P$都有使用半空间的交集或极限点的凸包的对偶表示$P=\{x:s_i^Tx\leqslant r_i,\forall\, i\}=\{x:x=\sum{a_jv_j,\sum{a_j}=1,a_j\geqslant 0,\forall\, j}\}$
  \end{proposition}
\end{frame}

\begin{frame}{多面体的半空间表示}
    基于以上命题，如果考虑来自$\mathbb{R}^d$中的有限个超平面
  \begin{equation}
    \mathcal{H}_i=\{x\in\mathbb{R}^d:h_ix=k_i\},i\in\mathcal{I}
  \end{equation}
  其中，$\mathcal{I}\triangleq \{1,\cdots,N\},(h_i,k_i)\in\mathbb{R}^{1\times d}\times\mathbb{R}$\par
  每一个超平面都会把空间分成两个不相交的区域，分别记作：
  \begin{subequations}
    \begin{gather}
    \mathcal{R}^+_i=\{x\in\mathbb{R}^d:h_ix\leqslant k_i\} \\
    \mathcal{R}^-_i=\{x\in\mathbb{R}^d:-h_ix\leqslant -k_i\}
  \end{gather}
  \end{subequations}\par
  则多面体$P$可写成这些半空间的的有界交集，另可得到其补集
  \begin{equation}
      P=\bigcap_{i\in\mathcal{I}}{\mathcal{R}^{-}_i}
  \end{equation}
  \begin{equation}
        \mathcal{C}(P)=\bigcup_{i\in\mathcal{I}}{\mathcal{R}^{+}_i} \label{6}
  \end{equation}
\end{frame}

\begin{frame}{使用二元变量表示区域}
  为了便于描述上述区域，可以引入二元变量$(\alpha_1,\alpha_N)\in\{0,1\}^N$，使用混合整数技术进行表示.\par
  \begin{subequations}
    \begin{gather}
      h_ix\leqslant k_i+M\alpha_i,i\in\mathcal{I} \label{7.a}\\
      \sum_{i\in\mathcal{I}}{\alpha_i}\leqslant N-1
    \end{gather}
  \end{subequations}
  \begin{example}
    半空间$R_x^+$可以由(\ref{7.a})和以下二元变量表示：
    \begin{equation}
      (\alpha_1,\cdots,\alpha_N)=\left(1\cdots 1\;\underbrace{0}_{i}\; 1\cdots 1\right)
    \end{equation}
  \end{example}
\end{frame}

%-----------------------------------------------
\section{MIP在运动规划中的应用}
%-----------------------------------------------



%------------------------------------------------
\section{多智能体系统中的MIP\\连接性维护和编队控制}
%------------------------------------------------





%------------------------------------------------
\section{控制架构}
%------------------------------------------------




%------------------------------------------------
\section{使用MIP的控制软件}
%------------------------------------------------





%------------------------------------------------
\section{MIP的替代方案}
%------------------------------------------------


\begin{frame}
\frametitle{Blocks of Highlighted Text}
\begin{block}{Block Title}
This is the block environment. The quick brown fox jumps over the lazy dog. The quick brown fox jumps over the lazy dog. The quick brown fox jumps over the lazy dog.
\end{block}

\begin{exampleblock}{Block Title}
This is the exampleblock environment. The quick brown fox jumps over the lazy dog. The quick brown fox jumps over the lazy dog.
\end{exampleblock}

\begin{alertblock}{Block Title}
This is the alertblock environment. The quick brown fox jumps over the lazy dog. The quick brown fox jumps over the lazy dog.
\end{alertblock}
\end{frame}

%------------------------------------------------

\begin{frame}
\frametitle{列表环境}

计数列表环境
\begin{enumerate}
\item 这是一个计数列表环境.
\item 这是一个计数列表环境.
\item 这是一个计数列表环境.
\end{enumerate}

\vspace{2ex}
不计数列表环境
\begin{itemize}[<+-| alert@+>]
\item 这是一个不计数列表环境.
\item 这是一个不计数列表环境.
\item 这是一个不计数列表环境.
\end{itemize}

\end{frame}

%------------------------------------------------

\begin{frame}
\frametitle{左右分栏}  % Multiple Columns
\begin{columns}[t] % The "c" option specifies centered vertical alignment while the "t" option is used for top vertical alignment

\column{.45\textwidth} % Left column and width
\textbf{Heading}
\begin{enumerate}
\item Statement
\item Explanation
\item Example
\end{enumerate}

\column{.5\textwidth} % Right column and width
The quick brown fox jumps over the lazy dog. The quick brown fox jumps over the lazy dog. The quick brown fox jumps over the lazy dog. The quick brown fox jumps over the lazy dog.
\end{columns}
\end{frame}


\begin{frame}
\frametitle{定理环境}
\begin{definition} \upshape\rmfamily
This is a definition environment. 这是一个定义环境.
\end{definition}


\begin{lemma} \upshape\rmfamily
This is a lemma environment. 这是一个引理环境.
\end{lemma}

\begin{proposition} \upshape\rmfamily
This is a proposition environment. 这是一个命题环境.
\end{proposition}

\begin{theorem}[Mass--energy] \upshape\rmfamily
This is a theorem environment. 这是一个定理环境.
\end{theorem}

\begin{proof}
  This is a proof environment. 这是一个证明环境.
\end{proof}

\end{frame}

%------------------------------------------------


\begin{frame}
\frametitle{定理示例}

\begin{theorem}[Lax-Milgram Lemma] \upshape
Let $X$ be a Hilbert space, let $a(\cdot, \cdot)$ : $X \times X \rightarrow \mathbb{R}$ be a continuous and coercive bilinear form, and let $F : X \rightarrow \mathbb{R}$ be a linear functional in $X^{\prime}$. Then the variational problem:
\begin{equation}
  \alert{
  \left\{\begin{aligned}
  &\text {Find } u \in X \text { such that } \\
  &a(u, v)=F(v), \forall v \in X
  \end{aligned} \right. }
\end{equation}
has a unique solution. Moreover, we have
\begin{equation}
  \alert{ \|u\| \leq \frac{1}{\alpha}\|F\|_{X^{\prime}}  }
\end{equation}
\end{theorem}

\end{frame}

%------------------------------------------------

\begin{frame}[fragile] % Need to use the fragile option when verbatim is used in the slide
\frametitle{Verbatim}
\begin{example}[Theorem Slide Code]
\begin{verbatim}
\begin{frame}
\frametitle{Theorem}
\begin{theorem}[Mass--energy equivalence]
$E = mc^2$
\end{theorem}
\end{frame}\end{verbatim}
\end{example}

\begin{table}
\caption{这是一个三线表.}
\begin{tabular}{l l l}
\toprule
\textbf{Treatments} & \textbf{Response 1} & \textbf{Response 2}\\
\midrule
Treatment 1 & 0.0003262 & 0.562 \\
Treatment 2 & 0.0015681 & 0.910 \\
\bottomrule
\end{tabular}
\end{table}

\end{frame}


%------------------------------------------------

\begin{frame}
\frametitle{表格环境}
本文定义了新的可变长度左中右 (LCR) 格式, LCR 三个格式会根据表格宽度的设定自行控制宽度, 且其宽度相等, 方便设置和页面相同宽度的表格. 本文还定义了 P\{\} 格式可以设定某一列宽度 (如 P\{1cm\} 控制某一列的宽度为 1cm) 并居中.
\begin{table}[!htp]
\centering
% PLCR已经定义
\caption{某校学生身高体重样本.}
\label{tab2:heightweight}
\begin{tabularx}{0.9\textwidth}{lCCC}
   \toprule
	序号&年龄&身高&体重\\
	\midrule
	1&14&156&42\\
	2&16&158&45\\
	3&14&162&48\\
	4&15&163&50\\
    \cmidrule{2-4}
	平均&15&159.75&46.25\\
	\bottomrule
\end{tabularx}
\end{table}

\end{frame}


%------------------------------------------------

\begin{frame}
\frametitle{表格示例}
\begin{table}[htp!]
\centering
\renewcommand\arraystretch{1.0} %定义表格高度
% PLCR前面已经定义
\caption{表格的描述.}
\label{tab3:NumError}
\begin{tabularx}{0.9\textwidth}{|P{1cm}|C|C|C|C|}
\Xhline{2\arrayrulewidth}
N  & A       & B    & C  & D   \\
\Xhline{2\arrayrulewidth}
1  & 9.20E-05 & 9.90E-05 & 1.00E-06 & 8.00E-06  \\
2  & 9.80E-05 & 8.00E-05 & 7.00E-06 & 1.40E-05  \\
3  & 4.00E-06 & 8.10E-05 & 8.80E-05 & 2.00E-05 \\
4  & 8.50E-05 & 8.70E-05 & 1.90E-05 & 2.10E-05 \\
5 & 8.60E-05 & 9.30E-05 & 2.50E-05 & 2.00E-06  \\
6 & 1.70E-05 & 2.40E-05 & 7.60E-05 & 8.30E-05  \\
7 & 2.30E-05 & 5.00E-06 & 8.20E-05 & 8.90E-05 \\
8 & 7.90E-05 & 6.00E-06 & 1.30E-05 & 9.50E-05  \\
9 & 1.00E-05 & 1.20E-05 & 9.40E-05 & 9.60E-05 \\
\Xhline{2\arrayrulewidth}
\end{tabularx}
\end{table}

\end{frame}


\begin{frame}[fragile] % Need to use the fragile option when verbatim is used in the slide
\frametitle{Citation}
An example of the \verb|\cite| command to cite within the presentation:\\~

This statement requires citation \cite{Smith2012}. \\~

文献引用示例 \cite{LiLiu1997}, 可以修改引用文献样式.
\end{frame}

%------------------------------------------------

\begin{frame}
\frametitle{References}
\footnotesize{
\begin{thebibliography}{99} % Beamer does not support BibTeX so references must be inserted manually as below
\bibitem[Smith, 2012]{Smith2012} John Smith, Title of the publication, \emph{Journal Name}, 12(3):45--678, 2012.
\bibitem[李荣华, 1997]{LiLiu1997} 李荣华, 刘播. 微分方程数值解法. 东南大学出版社, 1997.
\end{thebibliography}
}
\end{frame}


%------------------------------------------------
%\setbeamertemplate{background canvas}[vertical shading][bottom=white,top=structure.fg!25]
%\setbeamertemplate{headline}{}
%\begin{frame}
%\sffamily
%\Huge{\centerline{The End}}
%\end{frame}

%----------------------------------------------------------------------------------------
%\setbeamertemplate{background canvas}[vertical shading][bottom=white,top=structure.fg!25]
%\setbeamertemplate{headline}{}
%\begin{frame}
%\sffamily
%\begin{center}
%\HUGE{\textcolor[RGB]{165,3,3}{Thank~you!}}
%\end{center}
%\end{frame}


\begin{frame}[standout]
  Questions?
\end{frame}


\end{document}


